\documentclass[a4paper,11pt]{article}
\usepackage[T2A]{fontenc}
\usepackage[utf8]{inputenc}
\usepackage[russian]{babel}
\usepackage{amsmath,amssymb,mathtools}
\usepackage{geometry}
\usepackage{microtype}
\usepackage{tikz}
\usetikzlibrary{arrows.meta,positioning,shapes.geometric,calc}
\usepackage{pgfplots}
\pgfplotsset{compat=1.18}
\usepackage{booktabs,tabularx,array}
\usepackage{enumitem}
\usepackage{fancyhdr}
\usepackage{xcolor}
\usepackage{setspace}
\usepackage{hyperref}
\usepackage{needspace}
\usepackage{icomma}
\geometry{left=25mm,right=25mm,top=23mm,bottom=24mm}
\setstretch{1.20}
\definecolor{accent}{HTML}{008080}
\definecolor{darkaccent}{HTML}{005B5C}
\definecolor{maintext}{HTML}{333333}
\definecolor{muted}{HTML}{666666}
\color{maintext}
\hypersetup{colorlinks=true,linkcolor=darkaccent}
\setlength{\parindent}{0pt}
\setlength{\parskip}{5.5pt}
\setlist[itemize]{itemsep=2pt,topsep=2pt,leftmargin=6mm}
\setlist[enumerate]{itemsep=5pt,topsep=4pt,leftmargin=7mm}
\pagestyle{fancy}
\fancyhf{}
\fancyhead[L]{\footnotesize\color{muted}Логарифмические функции}
\fancyhead[R]{\footnotesize\color{muted}Денис Ермолаев}
\fancyfoot[C]{\footnotesize\thepage}
\renewcommand{\headrulewidth}{0.30pt}
\renewcommand{\headrule}{\hbox to\headwidth{\color{accent!60}\leaders\hrule height \headrulewidth\hfill}}
\renewcommand{\footrulewidth}{0pt}
\setlength{\headheight}{22pt}
\newcommand{\sectionrule}{\par\vspace{2pt}{\color{accent!65}\hrule height 0.42pt}\vspace{7pt}}
\newcommand{\topic}[1]{\needspace{7\baselineskip}\par\vspace{11pt}{\large\color{darkaccent}#1}\sectionrule}
\newcommand{\mini}[1]{\par\vspace{5pt}{\color{darkaccent}\bfseries #1}\par}
\newcommand{\note}[1]{{\color{darkaccent}\bfseries #1}}
\begin{document}
\thispagestyle{empty}
\begin{center}
\vspace*{27mm}
{\small\color{darkaccent}\MakeUppercase{Математика}}\par
\vspace{18mm}
{\LARGE Чек-лист}\par
\vspace{12mm}
{\Large\color{accent} Графики логарифмических функций\\[2mm]
и метод узловых точек}\par
\vspace{14mm}
{\normalsize Подготовка к ЕГЭ по профильной математике}\par
\vspace{43mm}
{\large Денис Ермолаев}\par
\vspace{12mm}
{\normalsize 2 ноября 2026 г.}\par
\vfill
{\small Лёвин Артём Александрович\\[2pt]
эксклюзивно для Дениса Ермолаева}\par
\vspace{16mm}
\begin{tikzpicture}[x=1cm,y=1cm]
\draw[accent,line width=.78pt] (0,0) .. controls (1.35,.58) and (3.05,-.52) .. (4.75,0)
 .. controls (6.40,.57) and (8.05,-.54) .. (9.7,.08);
\fill[darkaccent] (4.75,0) circle (1.1pt);
\end{tikzpicture}
\end{center}
\clearpage

\topic{1. Что нужно знать о логарифме}
По определению логарифм \(\log_a u\) --- это показатель степени, в которую нужно возвести \(a\), чтобы получить \(u\):
\[
\boxed{\log_a u=c\quad\Longleftrightarrow\quad a^c=u.}
\]
Условия существования действительного логарифма:
\[
\boxed{a>0,\qquad a\ne1,\qquad u>0.}
\]
Особенно полезны равенства
\[
\log_a1=0,\qquad \log_a a=1,\qquad
\log_{a}a^t=t,\qquad \log_{a^m}a^n=\frac{n}{m}\quad(m\ne0).
\]
\mini{Пример. Основание и аргумент --- степени одного числа}
\[
\log_8 4=\log_{2^3}2^2=\frac23.
\]
\note{Проверка:} \(8^{2/3}=4\). Замена \(\log_a u=c\) на \(a^c=u\) помогает избежать путаницы с показателями.

\topic{2. Как читать запрос в условии}
\begin{itemize}
\item \(f(32)\) --- требуется \emph{значение функции}: подставьте \(x=32\).
\item \(f(x)=1\) --- требуется \emph{аргумент}: решите уравнение относительно \(x\).
\end{itemize}
\mini{Пример}
Если \(f(x)=\log_2x-3\), то
\[
f(32)=\log_232-3=5-3=2.
\]
Если для той же функции \(f(x)=1\), то
\[
\log_2x-3=1\ \Longrightarrow\ \log_2x=4\ \Longrightarrow\ x=16.
\]
Второй пример требует также проверки \(x>0\); найденное значение допустимо.

\topic{3. Два разных вида логарифмической функции}
\begin{center}
\begin{tabularx}{\linewidth}{@{}X X@{}}
\textcolor{darkaccent}{\(f(x)=\log_a x+b\)} & \textcolor{darkaccent}{\(g(x)=\log_a(x+b)\)}\\[4pt]
Число \(b\) стоит \textbf{после} логарифма: сдвиг графика вверх/вниз. ОДЗ: \(x>0\). & Число \(b\) находится \textbf{под} логарифмом: сдвиг по горизонтали. ОДЗ: \(x+b>0\).\\
\end{tabularx}
\end{center}
При подстановке координат \((x_0;y_0)\) важно сначала точно записать формулу функции и только затем заменить \(x\) на \(x_0\), \(f(x)\) на \(y_0\).

\topic{4. Метод узловых точек}
\textbf{Узловая точка} --- достоверная точка графика с целочисленными координатами. Удобно выбирать точки, у которых координаты малы по модулю, но главное условие --- \emph{получить независимые уравнения}.

Если в функции два неизвестных параметра, обычно берут две подходящие точки, подставляют их координаты и решают систему. Затем проверяют ограничения и вычисляют искомое значение функции.

\topic{5. Модель \(f(x)=\log_a x+b\)}
Пусть графику принадлежат точки \(A(1;-3)\) и \(B(4;-1)\). Найдём функцию и \(f(32)\).
\[
\begin{cases}
-3=\log_a1+b,\\
-1=\log_a4+b.
\end{cases}
\quad\Longrightarrow\quad
\begin{cases}
b=-3,\\
\log_a4=2.
\end{cases}
\]
\(\log_a4=2\) означает \(a^2=4\). Из двух алгебраических корней \(a=2\) и \(a=-2\) остаётся только \(a=2\), потому что основание логарифма положительно.
\[
\boxed{f(x)=\log_2 x-3},\qquad
f(32)=5-3=\boxed{2}.
\]
\mini{Проверка по исходным точкам}
\(f(1)=0-3=-3\), \(f(4)=2-3=-1\). Обе точки подходят.

\topic{6. Модель \(g(x)=\log_a(x+b)\)}
Пусть графику принадлежат \(A(-4;0)\) и \(B(-1;2)\). Найдём \(g(11)\).
\[
\begin{cases}
0=\log_a(-4+b),\\
2=\log_a(-1+b).
\end{cases}
\]
\(0=\log_a1\), следовательно \(-4+b=1\) и \(b=5\). Тогда
\[
2=\log_a4\ \Longrightarrow\ a^2=4\ \Longrightarrow\ a=2.
\]
\[
\boxed{g(x)=\log_2(x+5)},\qquad
 g(11)=\log_216=\boxed{4}.
\]
\note{ОДЗ:} \(x>-5\). Проверка: \(g(-4)=\log_21=0\), \(g(-1)=\log_24=2\).

\topic{7. Знак «минус» и отрицательная степень}
Если \(h(x)=-\log_2(x+2)\) и \(h(x)=-4\), сначала уберите внешний минус:
\[
-\log_2(x+2)=-4\ \Longrightarrow\
\log_2(x+2)=4\ \Longrightarrow\ x+2=16.
\]
Отсюда \(\boxed{x=14}\), и \(14>-2\).

\note{Важно:} \(-\log_a u=\log_a(1/u)\), но \(-\log_a u\ne\log_a(-u)\). Кроме того, \(a^{-1}=1/a\): это свойство относится к основанию степени, если \(-1\) стоит в показателе.

\topic{8. Типичные ошибки и самопроверка}
\begin{itemize}
\item \textbf{Перепутать запрос:} \(f(16)\) и \(f(x)=16\) требуют разных действий.
\item \textbf{Подменить аргумент:} для \(\log_a(x+b)\) под логарифмом стоит целиком \(x+b\).
\item \textbf{Забыть ОДЗ:} основание \(a>0\), \(a\ne1\), аргумент положителен.
\item \textbf{Принять лишнюю точку за узловую:} координаты надо считывать точно, а не приблизительно.
\item \textbf{Поделить уравнения на ноль:} допустимы лишь преобразования с ненулевым делителем; система должна содержать независимые сведения.
\item \textbf{Пытаться найти оба коэффициента} в \(k\log_a x\): здесь параметры могут оказаться неидентифицируемыми.
\end{itemize}
\mini{Короткая разминка из начала занятия}
Для вычислений полезно переносить числа в степенной вид:
\[
0{,}6=6\cdot10^{-1},\qquad
4200=4{,}2\cdot10^3=42\cdot10^2.
\]
В стандартной научной записи перед \(10^n\) стоит число, модуль которого от \(1\) до \(10\) (не включая \(10\)).

\vspace{5mm}
\begin{center}\color{darkaccent}\textit{Сначала восстановите вид функции, затем найдите параметры,\\ после этого выполняйте действие, указанное в вопросе.}\end{center}

\clearpage
\topic{9. Два разных сдвига графика}
Сравните графики \(y=\log_2x-2\) и \(y=\log_2(x+2)\). Первый определён при \(x>0\), второй --- при \(x>-2\).
\vspace{1mm}
\begin{center}
\begin{tikzpicture}
\begin{axis}[
width=.92\linewidth,height=7.2cm,
xmin=-3,xmax=9,ymin=-4,ymax=4,
axis lines=middle, axis line style={gray!75,->},
xtick={-2,0,1,2,4,8},ytick={-3,-2,-1,0,1,2,3},
minor tick num=0,ticklabel style={font=\scriptsize},
grid=major,grid style={gray!12},
legend style={at={(.98,.04)},anchor=south east,draw=none,fill=white,font=\footnotesize},
clip=false]
\addplot[accent,thick,domain=.12:9,samples=200] {ln(x)/ln(2)-2};
\addlegendentry{\(y=\log_2x-2\)}
\addplot[darkaccent,densely dashed,thick,domain=-1.88:9,samples=220] {ln(x+2)/ln(2)};
\addlegendentry{\(y=\log_2(x+2)\)}
\addplot[only marks,mark=*,mark size=2pt,color=accent] coordinates {(1,-2) (4,0)};
\addplot[only marks,mark=square*,mark size=2pt,color=darkaccent] coordinates {(-1,0) (2,2)};
\draw[dotted,gray!70] (axis cs:-2,-4)--(axis cs:-2,4);
\end{axis}
\end{tikzpicture}
\end{center}
\note{Наблюдение:} \(\log_a(x+b)=0\) при \(x+b=1\). Для \(\log_a x+b=0\) требуется решить \(\log_a x=-b\).

\topic{10. Когда точки не позволяют найти параметры}
Если функция задана как
\[
f(x)=k\log_a x,\qquad a>0,\quad a\ne1,
\]
то из всех её точек вообще нельзя по отдельности восстановить оба неизвестных числа \(k\) и \(a\):
\[
 k\log_a x=\frac{k}{\log_2 a}\,\log_2x.
\]
График определяет лишь \emph{отношение} \(k/\log_2 a\). Например, функции \(3\log_2x\) и \(6\log_4x\) совпадают.

\mini{Как извлечь пользу из точек}
Если \(A(2;3)\) и \(B(4;6)\) принадлежат графику \(k\log_a x\), то
\[
f(8)=k\log_a 2^3=3k\log_a2=3f(2)=\boxed{9}.
\]
Нужное значение однозначно, хотя \(a\) и \(k\) по отдельности неизвестны.

\clearpage
\thispagestyle{fancy}
\begin{center}
{\Large\color{darkaccent}Алгоритм: функция по узловым точкам}
\end{center}
\vspace{9mm}
\begin{center}
\tikzset{
startstop/.style={ellipse,draw=darkaccent!70,very thin,align=center,text width=5.7cm,minimum height=10mm,fill=white},
action/.style={rectangle,rounded corners=2pt,draw=accent!68,very thin,align=center,text width=8.1cm,minimum height=12mm,fill=white},
branch/.style={diamond,aspect=2.6,draw=darkaccent!75,very thin,align=center,text width=4.5cm,inner ysep=5pt,fill=white},
arr/.style={-{Latex[length=2.1mm]},line width=.7pt,color=darkaccent!78}}
\begin{tikzpicture}[node distance=10mm and 13mm, every node/.style={font=\small}]
\node[startstop] (s) {Начало: прочитайте вид функции\\и то, что требуется найти};
\node[action,below=of s] (t) {Выберите точные узловые точки;\\учтите ОДЗ и ограничения основания};
\node[action,below=of t] (u) {Подставьте \((x_i;y_i)\): составьте\\уравнения для неизвестных параметров};
\node[branch,below=10mm of u] (d) {Данных достаточно?\\Условия независимы?};
\node[action,below=13mm of d] (v) {Решите систему, проверьте \(a>0\), \(a\ne1\);\\восстановите функцию и проверьте точки};
\node[action,right=13mm of d,text width=4.2cm] (w) {Смените точки. Если параметры неразличимы, ищите нужное значение напрямую};
\node[action,below=11mm of v] (f) {Подставьте заданный \(x\) или\\решите уравнение \(f(x)=c\)};
\node[startstop,below=of f] (e) {Проверьте ответ и ОДЗ};
\draw[arr] (s)--(t);
\draw[arr] (t)--(u);
\draw[arr] (u)--(d);
\draw[arr] (d)--node[left,font=\footnotesize]{да}(v);
\draw[arr] (d.east)--node[above,font=\footnotesize]{нет}(w.west);
\draw[arr] (w.south) |- (f.east);
\draw[arr] (v)--(f);
\draw[arr] (f)--(e);
\end{tikzpicture}
\end{center}
\clearpage

\clearpage

\begin{center}
{\Large\color{darkaccent}Тренировочный блок}\par
\vspace{2mm}
{\small Путь А \(\cdot\) 10 заданий \(\cdot\) от базовых к содержательным}
\end{center}
\vspace{5mm}
В задачах с неизвестным основанием проверяйте ограничения для логарифма. Если требуется, укажите ОДЗ.

\mini{Базовый уровень}
\begin{enumerate}[label=\textcolor{darkaccent}{\arabic*.},start=1]
\item Вычислите \(\displaystyle \log_8 4\).
\item Вычислите \(\displaystyle \log_3\frac1{27}\).
\item Для \(f(x)=\log_2x-3\) найдите \(f(16)\).
\item Для \(f(x)=\log_2(x+3)\) найдите \(x\), если \(f(x)=3\).
\end{enumerate}
\mini{Средний уровень}
\begin{enumerate}[label=\textcolor{darkaccent}{\arabic*.},start=5]
\item Функция \(f(x)=\log_ax+b\) проходит через \(A(1;-2)\) и \(B(9;0)\). Найдите \(f(27)\).
\item Функция \(g(x)=\log_a(x+b)\) проходит через \(A(-1;0)\) и \(B(2;2)\). Найдите \(g(6)\).
\item Функция \(f(x)=\log_ax+b\) проходит через \(A(1;1)\) и \(B(8;4)\). Найдите \(x\), если \(f(x)=2\).
\end{enumerate}
\mini{Уровень выше среднего}
\begin{enumerate}[label=\textcolor{darkaccent}{\arabic*.},start=8]
\item Функция \(g(x)=\log_a(x+b)\) проходит через \(A(0;0)\) и \(B(3;2)\). Найдите \(x\), если \(g(x)=-3\).
\item Функция \(f(x)=k\log_2x\) проходит через \(A(8;-6)\). Вычислите \(f(4)\).
\item График \(f(x)=k\log_ax\) проходит через \(A(2;3)\) и \(B(4;6)\). Найдите \(f(8)\). Можно ли по этим данным однозначно определить \(a\) и \(k\) по отдельности?
\end{enumerate}
\clearpage

\begin{center}
{\Large\color{darkaccent}Ответы}\par
\vspace{3mm}
{\small К тренировочному блоку из 10 заданий}
\end{center}
\vspace{13mm}
\renewcommand{\arraystretch}{1.65}
\begin{center}
\begin{tabularx}{.92\linewidth}{@{}>{\centering\arraybackslash}p{20mm} X@{}}
\toprule
\textbf{№} & \textbf{Ответ}\\
\midrule
1 & \(\displaystyle \frac23\)\\
2 & \(-3\)\\
3 & \(1\)\\
4 & \(x=5\)\\
5 & \(1\)\\
6 & \(3\)\\
7 & \(x=2\)\\
8 & \(\displaystyle x=-\frac78\)\\
9 & \(-4\)\\
10 & \(f(8)=9\). Нет: \(a\) и \(k\) по отдельности не определяются.\\
\bottomrule
\end{tabularx}
\end{center}
\vfill
\begin{center}
{\small\color{muted}Лёвин Артём Александрович\\
эксклюзивно для Дениса Ермолаева}
\end{center}
\end{document}
